\documentclass[11pt]{article}

\usepackage{amsmath,amssymb,amsfonts,mathrsfs}
\usepackage{hyperref}
\usepackage{geometry}
\geometry{a4paper,margin=1in}

\title{Wigner--In\"on\"u Contractions}
\author{}
\date{}

\begin{document}

\maketitle


The Wigner--In\"on\"u contraction is a systematic procedure that produces one Lie
algebra as a singular limit of another.  Originally introduced by Wigner and
In\"on\"u~\cite{InonuWigner1953}, the method explains how nonrelativistic or
flat--space symmetry groups arise algebraically from relativistic or curved
ones.  The contraction operates entirely at the level of Lie algebras and does
not in general correspond to an isomorphism; it is a degeneration that reduces
the structure constants and typically turns a semisimple algebra into a
non-semisimple one.

\section{Definition}

Let $\mathfrak{g}$ be a Lie algebra with basis $\{X_i\}$ and commutation
relations
\begin{equation}
   [X_i,X_j] = f_{ij}{}^{k} X_k ,
\end{equation}
with constant structure coefficients $f_{ij}{}^{k}$.  A
Wigner--In\"on\"u contraction of $\mathfrak{g}$ is defined as follows:

\begin{enumerate}
\item Introduce a one--parameter family of invertible linear transformations
      $X_i \mapsto X_i(\varepsilon)$ for $\varepsilon \neq 0$.
\item Express the commutators in the $\varepsilon$--dependent basis:
\begin{equation}
    [X_i(\varepsilon),X_j(\varepsilon)]
        = f_{ij}{}^{k}(\varepsilon) \, X_k(\varepsilon).
\end{equation}
\item Take the limit $\varepsilon \to 0$, assuming it exists:
\begin{equation}
    \lim_{\varepsilon\to 0} f_{ij}{}^{k}(\varepsilon)
        = f^{(0)}_{ij}{}^{k}.
\end{equation}
\end{enumerate}

The algebra determined by $f^{(0)}_{ij}{}^{k}$ is called the
\emph{contracted Lie algebra} and is denoted $\mathfrak{g}_0$.

\section{Simple Wigner--In\"on\"u Contraction}

In the original formulation, $\mathfrak{g}$ is assumed to decompose as a vector
space into a direct sum
\begin{equation}
    \mathfrak{g} = \mathfrak{h} \oplus \mathfrak{p},
\end{equation}
where $\mathfrak{h}$ is a subalgebra.  One then rescales only the generators in
$\mathfrak{p}$:
\begin{equation}
   X_h' = X_h, \qquad
   X_p' = \varepsilon\, X_p,
\end{equation}
with $X_h \in \mathfrak{h}$ and $X_p \in \mathfrak{p}$.
Taking $\varepsilon \to 0$ yields the contracted algebra.

\section{Example: Poincar\'e \texorpdfstring{$\to$}{→} Galilei}

Let $J_i$ be spatial rotations, $K_i$ Lorentz boosts, $P_i$ spatial momenta,
and $H$ the energy operator of the Poincar\'e algebra.  The characteristic
relativistic commutators are
\begin{equation}
   [K_i,K_j] = - \frac{1}{c^2} \epsilon_{ijk} J_k,
   \qquad
   [K_i,P_j] = \frac{1}{c^2} \delta_{ij} H.
\end{equation}
Define rescaled generators
\begin{equation}
    K_i' = c\, K_i, \qquad H' = H.
\end{equation}
Then
\begin{equation}
   [K_i',K_j'] = - c^2 \frac{1}{c^2} \epsilon_{ijk} J_k
      \xrightarrow[c\to\infty]{} 0, 
   \qquad
   [K_i',P_j] = c \frac{1}{c^2} \delta_{ij} H
      \xrightarrow[c\to\infty]{} 0 .
\end{equation}
The resulting commutation relations coincide with those of the Galilei algebra,
showing that the Galilean symmetry arises as the singular limit $c\to\infty$
of the Poincar\'e symmetry.

\section{Example: \texorpdfstring{$\mathfrak{so}(3) \to \mathfrak{e}(2)$}{so(3)→e(2)}}

Let $J_x,J_y,J_z$ generate $\mathfrak{so}(3)$.  Define
\begin{equation}
    P_x = \varepsilon J_x, \qquad
    P_y = \varepsilon J_y, \qquad
    J = J_z.
\end{equation}
Then
\begin{align}
    [P_x,P_y] &= \varepsilon^2 J_z \xrightarrow{\varepsilon\to 0} 0, \\
    [J,P_x] &= P_y, \qquad [J,P_y] = - P_x.
\end{align}
This is precisely the Lie algebra $\mathfrak{e}(2)$ of the Euclidean group
$E(2)$ in two dimensions.  Geometrically, this corresponds to the limit in
which a sphere of radius $R = 1/\varepsilon$ becomes flat as $R\to\infty$.

\section{Generalized In\"on\"u--Wigner Contractions}

The simple contraction may be generalized by assigning distinct scaling
exponents to different generators:
\begin{equation}
    X_i' = \varepsilon^{n_i} X_i ,
\end{equation}
with nonnegative integers $n_i$.  This includes all graded contractions and
produces a wide variety of non-isomorphic limit algebras.  Such generalized
contractions underlie the classification of kinematical Lie algebras by Bacry
and L\'evy-Leblond.

\section{Physical Interpretation}

Contractions correspond to well-defined physical limits: the
nonrelativistic limit $c\to\infty$, flat-space limits $R\to\infty$ of curved
spacetimes, ultra-relativistic limits leading to Carroll algebras, and
weak-coupling or large-mass limits in field theory.  In all cases, the
contraction captures the algebraic residue of a theory after certain dynamical
effects have been suppressed.



\end{document}

